\chapter{注意力机制基础}

\section{引言}

注意力机制（Attention Mechanism）是现代深度学习的核心技术之一，尤其在自然语言处理领域取得了巨大成功。本节回顾标准注意力机制的基本原理，为后续理解线性注意力奠定基础。

\section{标准注意力公式}

注意力机制的核心思想是：给定一组查询（Query）和一组键值对（Key-Value），计算查询与每个键的相似度，并以此作为权重对值进行加权求和。

对于序列长度为 $N$、隐藏维度为 $d$ 的输入，标准注意力的计算过程如下：

\begin{equation}
    \mathbf{Q} = \mathbf{X}\mathbf{W}_Q, \quad
    \mathbf{K} = \mathbf{X}\mathbf{W}_K, \quad
    \mathbf{V} = \mathbf{X}\mathbf{W}_V
\end{equation}

其中 $\mathbf{X} \in \mathbb{R}^{N \times d}$ 是输入序列，$\mathbf{W}_Q, \mathbf{W}_K, \mathbf{W}_V \in \mathbb{R}^{d \times d}$ 是可学习的投影矩阵。

注意力输出为：

\begin{equation}
    \text{Attention}(\mathbf{Q}, \mathbf{K}, \mathbf{V}) = \text{softmax}\left(\frac{\mathbf{Q}\mathbf{K}^\top}{\sqrt{d}}\right)\mathbf{V}
\end{equation}

其中 $\text{softmax}$ 操作按行进行，$\sqrt{d}$ 是缩放因子。

\section{多头注意力}

Transformer 模型使用多头注意力（Multi-Head Attention），将查询、键、值分别投影到多个子空间中独立计算注意力，然后将结果拼接：

\begin{equation}
    \text{MultiHead}(\mathbf{Q}, \mathbf{K}, \mathbf{V}) = \text{Concat}(\text{head}_1, \ldots, \text{head}_h)\mathbf{W}_O
\end{equation}

其中：

\begin{equation}
    \text{head}_i = \text{Attention}(\mathbf{Q}\mathbf{W}_Q^i, \mathbf{K}\mathbf{W}_K^i, \mathbf{V}\mathbf{W}_V^i)
\end{equation}

$h$ 是头数，$\mathbf{W}_O$ 是输出投影矩阵。

\section{计算复杂度分析}

标准注意力的计算复杂度主要集中在 $\mathbf{Q}\mathbf{K}^\top$ 这一矩阵乘法上。对于序列长度 $N$ 和隐藏维度 $d$：

\begin{itemize}
    \item 计算 $\mathbf{Q}\mathbf{K}^\top$：$\mathcal{O}(N^2 d)$
    \item 计算 $\text{softmax}$：$\mathcal{O}(N^2)$
    \item 计算与 $\mathbf{V}$ 的乘积：$\mathcal{O}(N^2 d)$
    \item 总体时间复杂度：$\mathcal{O}(N^2 d)$
    \item 总体空间复杂度：$\mathcal{O}(N^2)$（需要存储 $N \times N$ 的注意力矩阵）
\end{itemize}

当序列长度 $N$ 增大时，二次复杂度成为显著瓶颈。例如，当 $N = 10^5$ 时，注意力矩阵需要存储 $10^{10}$ 个元素，这在实际中是不可行的。

\section{长序列挑战}

在实际应用中，许多场景需要处理超长序列：

\begin{itemize}
    \item 文档级文本理解：数万 token
    \item 高分辨率图像处理：$1024 \times 1024$ 图像展平后超过 $10^6$ token
    \item 基因组序列分析：百万级碱基对
    \item 音频信号处理：长时间音频序列
\end{itemize}

标准注意力机制在这些场景下计算和内存开销过大，这推动了高效注意力机制的研究，其中线性注意力是最重要的方向之一。
